\documentclass[a4paper,12pt]{paper}

\usepackage{mycommands}
\usepackage{amsmath}

\title{Problem Set 2}
\author{David Al-Attar \\ Michaelmas Term, \the\year}

\begin{document}

\maketitle

\begin{enumerate}

\item Consider the regularised least squares solution to a linear(ised) inverse problem obtained by minimising the function
\begin{equation*}
    J(\mathbf{m})=\frac{1}{2}(\mathbf{A}\mathbf{m}-\mathbf{d})^{T}\mathbf{C}^{-1}(\mathbf{A}\mathbf{m}-\mathbf{d})+\frac{\lambda}{2}\mathbf{m}^{T}\mathbf{B}\mathbf{m}.
\end{equation*}
Suppose that the data takes the form $\mathbf{d}=\mathbf{A}\mathbf{m}_{in}$ with $\mathbf{m}_{in}$ a specified input model. Show that the resulting least squares solution, $\mathbf{m}_{out}$, takes the form
\begin{equation*}
    \mathbf{m}_{out}=\mathbf{R}\mathbf{m}_{in},
\end{equation*}
where $\mathbf{R}$ is a matrix to be determined. Discuss the significance of this so-called resolution matrix, explaining why, in particular, it would be desirable for its value to be close to the identity matrix.

\item Consider the deformation of an isotropic elastic body due to an applied surface load, $\sigma$. The relevant linearised equations of motion are
\begin{gather*}
    -\frac{\partial T_{ij}}{\partial x_{j}}=0; \\
    T_{ij}=\lambda\frac{\partial u_{k}}{\partial x_{k}}\delta_{ij}+\mu\left(\frac{\partial u_{i}}{\partial x_{j}}+\frac{\partial u_{j}}{\partial x_{i}}\right),
\end{gather*}
within the reference body $M$, and subject to boundary conditions
\begin{equation*}
    T_{ij}\hat{n}_{j}=\sigma g_{i}
\end{equation*}
on its boundary, $\partial M$. Here $g_{i}$ is the acceleration due to gravity prior to the deformation. Note that self-gravitation is being ignored within these equations.

Suppose that observations, $\mathbf{u}_{i}^{obs}$, of the displacement vector have been made at surface locations $\mathbf{x}_i$ for $i=1,...,m$. An inverse problem could then be formulated to estimate the Lamé parameters, $\lambda$ and $\mu$, as well as the load, $\sigma$. To this end, we can seek to minimise the misfit
\begin{equation*}
    J=\frac{1}{2}\sum_{i=1}^{m}\frac{1}{\sigma_{i}^2}||\mathbf{u}(\mathbf{x}_{i})-\mathbf{u}_{i}^{obs}||^{2}
\end{equation*}
where $\sigma_{i}$ is the standard deviation for the $i$-th observation. Adapt the approach within Lecture 19 to derive sensitivity kernels for $J$ with respect to $\lambda$, $\mu$, and $\sigma$, showing how they can be calculated through one solution of the forward problem and one of an appropriate adjoint problem.

\item Within Lecture 19 the idea of a cross correlation travel time was introduced. This is based on finding the shift, $\overline{\tau}$, that maximises the cross correlation
\begin{equation*}
    C(\tau)=\int_{-\infty}^{\infty}s^{obs}(t-\tau)s(\mathbf{x}_{s},t)\dd t,
\end{equation*}
between an observed waveform, $s^{obs}$, at a surface location $\mathbf{x}_{s}$ and a synthetic one, $s$, both having been suitably windowed. By perturbing the optimality condition
\begin{equation*}
    \frac{dC}{d\tau}(\overline{\tau})=0
\end{equation*}
with respect to $s$, show that to first-order accuracy
\begin{equation*}
    \delta\overline{\tau}=\frac{\int_{-\infty}^{\infty}\dot{s}^{obs}(t-\overline{\tau})\delta s(\mathbf{x}_{s},t)\dd t}{\int_{-\infty}^{\infty}\ddot{s}^{obs}(t-\overline{\tau})s(\mathbf{x}_{s},t)\dd t},
\end{equation*}
where dots denote time derivatives. From this result, it is a simple matter to obtain adjoint tractions required for the calculation of banana doughnut kernels.

\item Recall that the Lagrangian for a self-gravitating elastic body is given by
\begin{equation*}
    \mathcal{L}(\mathbf{x},t,\boldsymbol{\varphi},\mathbf{v},\mathbf{F})=\frac{1}{2}\rho(\mathbf{x})||\mathbf{v}||^{2}-W(\mathbf{x},t,\mathbf{F})-\frac{1}{2}\rho(\mathbf{x})\zeta(\mathbf{x},t),
\end{equation*}
where $\zeta$ is the referential gravitational potential
\begin{equation*}
    \zeta(\mathbf{x},t)=-G\int_{M}\frac{\rho(\mathbf{x}')}{\{[\varphi_{i}(\mathbf{x},t)-\varphi_{i}(\mathbf{x}',t)][\varphi_{i}(\mathbf{x},t)-\varphi_{i}(\mathbf{x}',t)]\}^{\frac{1}{2}}}\dd^{3}\mathbf{x}'.
\end{equation*}
By calculating the appropriate functional derivative, show that
\begin{equation*}
    \frac{\delta\mathcal{L}}{\delta\varphi_{i}}=\rho\gamma_{i}
\end{equation*}
where $\gamma_{i}$ is the referential gravitational acceleration.

\item Consider a planet at equilibrium occupying the volume $M$ and with density $\rho$. Its gravitational potential, $\phi$, is a solution of Poisson's equation
\begin{equation*}
    \nabla^{2}\phi=4\pi G\rho,
\end{equation*}
where it is understood that $\rho$ vanishes outside of $M$. Across the surface $\partial M$ we then have continuity conditions on $\phi$ and its normal derivative, and finally require that $\phi$ tend to zero at infinity. Suppose that the planet is only slightly aspherical, with boundary
\begin{equation*}
    r=b+sh_{1}(\theta,\varphi)+\cdot\cdot\cdot,
\end{equation*}
in spherical polar co-ordinates, and its density decomposed as
\begin{equation*}
    \rho=\rho_{0}+s \rho_{1}+\cdot\cdot\cdot,
\end{equation*}
with $\rho_{0}$ spherically symmetric. Here $s$ is a perturbation parameter, and the dots denote higher-order terms that are to be neglected. Similarly writing the potential in the form $\phi=\phi_{0}+s \phi_{1}+\cdot\cdot\cdot$ show that across the reference surface $r=b$ we require continuity of both $\phi_{1}$ and
\begin{equation*}
    \frac{\partial\phi_{1}}{\partial r}+4\pi G\rho_{0}h_{1}.
\end{equation*}

\item Using results from Lecture 22 along with residue calculus (which you don't need to know for this course), it can be shown that the time-domain solution of the elastodynamic equations in a non-rotating earth model can be expressed as the following eigenfunction expansion
\begin{equation*}
    \mathbf{u}(\mathbf{x},t)=\sum_{km}\left\{\int_{0}^{t}\frac{\sin[\omega_{k}(t-t')]}{\omega_{k}}\int_{M}\overline{S}_{ij}(\mathbf{x}',t')\frac{\partial|km\rangle^{*}}{\partial x'_{j}}(\mathbf{x}')\dd^{3}\mathbf{x}'\dd t'\right\}|km\rangle(\mathbf{x}).
\end{equation*}
Simplify this expression in the case of a point source for which
\begin{equation*}
    \overline{S}_{ij}(\mathbf{x},t)=M_{ij}\delta(\mathbf{x}-\mathbf{x}_{s})H(t-t_{s}),
\end{equation*}
with $M_{ij}$ known as the moment tensor, $\mathbf{x}_{s}$ the source location, and $t_{s}$ the source time. Describe briefly and qualitatively how an inverse problem might be formulated to estimate the ten source parameters $(M_{ij},\mathbf{x}_{s},t_{s})$ from recorded seismograms.

\item Consider the eigenvalue problem for a rotating earth model
\begin{equation*}
    -\omega^{2}\langle \mathbf{w}|\mathbf{P}|\mathbf{s}\rangle+i\omega\langle \mathbf{w}|\mathbf{W}|\mathbf{s}\rangle+\langle \mathbf{w}|\mathbf{H}|\mathbf{s}\rangle=0
\end{equation*}
where $\mathbf{s}$ is the eigenfunction, $\omega$ the eigenvalue, and $\mathbf{w}$ an arbitrary test function. Assuming the eigenfunction normalisation condition
\begin{equation*}
    \langle \mathbf{s}|\mathbf{P}|\mathbf{s}\rangle=1,
\end{equation*}
show that the eigenfrequency can be written
\begin{equation*}
    \omega=\frac{1}{2}i\langle \mathbf{s}|\mathbf{W}|\mathbf{s}\rangle\pm\sqrt{\langle \mathbf{s}|\mathbf{H}|\mathbf{s}\rangle-\frac{1}{4}\langle \mathbf{s}|\mathbf{W}|\mathbf{s}\rangle^{2}}.
\end{equation*}
Use this result to discuss qualitatively the stability of a rotating earth model, including, in particular, the role of the Coriolis and centrifugal forces.

\end{enumerate}

\end{document}